% GENERATED FILE. Edit canonical lessons, then run npm run generate:books.
% canonical-source-hash: fnv1a64:8ec721f317505a9d
% canonical-lessons: LA-C56-sextus, LA-C56-septimus, LA-C56-octavus, LA-C56-nonus, LA-C56-decimus

\chapter{Sextus ad Decimum}
\label{ch:la-ordinals-to-tenth}
\hlchaptermodality{56}

\begin{chapteropening}
\textbf{By the end of this chapter:} \emph{I can count in order all the way to tenth, and say why noon and siesta are ordinal numbers wearing a disguise.}

The tenth ordinal, and the whole series said in order from primus to decimus.
\end{chapteropening}

\section[sextus, sexta, sextum]{sextus — the sixth, and the hour that became a nap}

\label{lesson:LA-C56-sextus}



\begin{quote}
Recalls open the chapter, at the four measured distances.

\begin{itemize}
\raggedright
  \item \emph{Recall:} say \emph{quīntus} — \textbf{R1}, one lesson back, before it decays
  \item \emph{Recall:} say \emph{prīmus} — \textbf{R2}, the first real retrieval, five to fifteen lessons back
  \item \emph{Recall:} say \emph{domus} — \textbf{R3}, durable, twenty to sixty lessons back
  \item \emph{Recall:} say \emph{ubi} — \textbf{R3}, a second word from the same distance
  \item \emph{Recall:} say \emph{sex, septem, octō, novem, decem} — \textbf{R4}, recognition at distance, eighty lessons back or more
\end{itemize}
\end{quote}

\begin{cousinweb}
\textbf{sextus, sexta, sextum} — \textquotedblleft{}\textbf{sixth}.\textquotedblright{}

\emph{Sex} is \textbf{six}; \emph{sextus} is the same word with the ending. And \emph{Sextīlis}, the other renamed month, is what \emph{Quīntīlis} was: an ordinal doing duty as a name, later handed over to Augustus.

The better story is the hour. A Roman day was cut into twelve hours running from sunrise to sunset, so \emph{prīma hōra} began at dawn and the \textbf{sixth} hour landed on the middle of the day — which is why \emph{sexta hōra} meant noon and the heat of it. That hour became a rest, the rest kept the hour's name, and the name wore down through Spanish into \textbf{siesta}. A word for a nap that is, underneath, an ordinal number.

English took the plainer ones: a \textbf{sextet}, a \textbf{sextant} (a sixth of a circle), and \emph{sexagēsimus} behind \textbf{sexagesimal}, counting by sixties.
\end{cousinweb}

\subsection*{Your turn}
Say these aloud:

\begin{itemize}
\raggedright
  \item \emph{sextus, sexta, sextum}
  \item the join — \emph{sex $\to$ sextus}
  \item \emph{sexta hōra} — the sixth hour, the middle of the day
\end{itemize}

\begin{tcolorbox}[breakable,colback=teal!4,colframe=teal!35!black,title={Before you move on}]
Which hour of a Roman day was \emph{sexta hōra}, and where did it fall? (\textbf{The sixth}, counted from sunrise — so \textbf{the middle of the day}.) What everyday Spanish word came out of it? (\emph{\textbf{Siesta}}.) And which month was \emph{Sextīlis} renamed after? (\textbf{Augustus}.)
\end{tcolorbox}
\section[septimus, septima, septimum]{septimus — the seventh, and the week that is a seven}

\label{lesson:LA-C56-septimus}



\begin{quote}
Recalls first, then the seventh.

\begin{itemize}
\raggedright
  \item \emph{Recall:} say \emph{sextus} — \textbf{R1}, one lesson back, before it decays
  \item \emph{Recall:} say \emph{secundus} — \textbf{R2}, the first real retrieval, five to fifteen lessons back
  \item \emph{Recall:} say \emph{porta} — \textbf{R3}, durable, twenty to sixty lessons back
  \item \emph{Recall:} say \emph{cūr} — \textbf{R3}, a second word from the same distance
  \item \emph{Recall:} say \emph{quaesō} — \textbf{R4}, recognition at distance, eighty lessons back or more
  \item \emph{Recall:} say \emph{ignōsce mihi} — \textbf{R4}, a second word from the same distance
\end{itemize}
\end{quote}

\begin{cousinweb}
\textbf{septimus, septima, septimum} — \textquotedblleft{}\textbf{seventh}.\textquotedblright{}

\emph{Septem} is \textbf{seven}. Put the ending on and the vowel in the middle goes quiet: \emph{septem} $\to$ \emph{septimus}, with \emph{e} worn down to \emph{i} because nothing is leaning on it. Latin does that to unstressed middles constantly, and once you have heard it here you will hear it everywhere.

\textbf{September} is this number in a name, and the calendar lesson explained why it sits in ninth place: two months were pushed in at the front and the names never moved.

The prettier descendant is the week. \emph{Septimāna} was \textquotedblleft{}a set of seven\textquotedblright{} — seven days — and it went into Spanish as \emph{semana} and French as \emph{semaine}. So in two of Latin's daughters the ordinary word for a week is a number word, and the number is this one.
\end{cousinweb}

\subsection*{Your turn}
Say these aloud:

\begin{itemize}
\raggedright
  \item \emph{septimus, septima, septimum}
  \item the join — \emph{septem $\to$ septimus}
  \item \emph{septima hōra} — the seventh hour
\end{itemize}

\begin{tcolorbox}[breakable,colback=teal!4,colframe=teal!35!black,title={Before you move on}]
What happens to the vowel in the middle of \emph{septem} when the ending goes on? (\textbf{It wears down} — \emph{e} to \emph{i}.) Which month keeps this number, and why is it two places out? (\emph{\textbf{September}}; \textbf{two months were added at the front}.) And what is a \emph{semana}, taken apart? (\textbf{A set of seven} — \emph{septimāna}.)
\end{tcolorbox}
\section[octāvus, octāva, octāvum]{octāvus — the eighth, and the ending that breaks step}

\label{lesson:LA-C56-octavus}



\begin{quote}
Recalls first, then the eighth — the one that does not do what the last four did.

\begin{itemize}
\raggedright
  \item \emph{Recall:} say \emph{septimus} — \textbf{R1}, one lesson back, before it decays
  \item \emph{Recall:} say \emph{tertius} — \textbf{R2}, the first real retrieval, five to fifteen lessons back
  \item \emph{Recall:} say \emph{mūrus} — \textbf{R3}, durable, twenty to sixty lessons back
  \item \emph{Recall:} say \emph{quia} — \textbf{R3}, a second word from the same distance
  \item \emph{Recall:} say \emph{diēs Lūnae} — \textbf{R4}, recognition at distance, eighty lessons back or more
  \item \emph{Recall:} say \emph{diēs Sōlis} — \textbf{R4}, a second word from the same distance
\end{itemize}
\end{quote}

\begin{cousinweb}
\textbf{octāvus, octāva, octāvum} — \textquotedblleft{}\textbf{eighth}.\textquotedblright{}

\emph{Octō} is \textbf{eight}. By now you would expect \emph{octus}, and it is not that. This one takes \textbf{-āvus}, and it is the only ordinal in the ten that does.

Worth stopping on, because it shows what a pattern in a language actually is. Four ordinals in a row behaved alike, so the shape looked settled; the eighth breaks step, and then the ninth breaks it differently, and the tenth goes back to behaving. A learner who was told a rule and then met these would think the rule was broken. A learner who was shown the shape and then shown where it bends has learned the language instead.

\textbf{October} is this number in a month's name. And an \textbf{octave} in music is the eighth note, the one at which the names begin again — \emph{do} to \emph{do} — so the interval is named for the count that closes it.
\end{cousinweb}

\subsection*{Your turn}
Say these aloud:

\begin{itemize}
\raggedright
  \item \emph{octāvus, octāva, octāvum}
  \item the ending that breaks step — not \emph{octus}
  \item the eight so far, in order
\end{itemize}

\begin{tcolorbox}[breakable,colback=teal!4,colframe=teal!35!black,title={Before you move on}]
What ending does \emph{octāvus} take, and how many of the ten take it? (\textbf{-āvus}, and \textbf{only this one}.) What is an octave, counted out? (\textbf{The eighth note} — where the names start again.) And what is \emph{October}, taken literally? (\textbf{The eighth month}.)
\end{tcolorbox}
\section[nōnus, nōna, nōnum]{nōnus — the ninth, and where noon comes from}

\label{lesson:LA-C56-nonus}



\begin{quote}
Recalls first, then the ninth — and the best word-history in the chapter.

\begin{itemize}
\raggedright
  \item \emph{Recall:} say \emph{octāvus} — \textbf{R1}, one lesson back, before it decays
  \item \emph{Recall:} say \emph{quārtus} — \textbf{R2}, the first real retrieval, five to fifteen lessons back
  \item \emph{Recall:} say \emph{hortus} — \textbf{R3}, durable, twenty to sixty lessons back
  \item \emph{Recall:} say \emph{meus} — \textbf{R3}, a second word from the same distance
  \item \emph{Recall:} say \emph{niger and albus} — \textbf{R4}, recognition at distance, eighty lessons back or more
  \item \emph{Recall:} say \emph{ruber and caeruleus} — \textbf{R4}, a second word from the same distance
\end{itemize}
\end{quote}

\begin{cousinweb}
\textbf{nōnus, nōna, nōnum} — \textquotedblleft{}\textbf{ninth}.\textquotedblright{}

\emph{Novem} is \textbf{nine}, and this time the word is not lengthened but squeezed: \emph{noven-} collapses in on itself and comes out as \emph{nōnus}, with a long \emph{ō} where the \emph{ove} used to be. Two ordinals running have now bent the shape in two different ways.

\textbf{November} is this number in a name.

And then \emph{nōna hōra}. Counting from sunrise, the ninth hour is the middle of the afternoon, and it was the hour of a fixed prayer. Over centuries that prayer crept earlier and earlier in the day, until it was being said at midday — and the \textbf{name went with it}. English \textbf{noon} is \emph{nōna}, still meaning \textquotedblleft{}the ninth\textquotedblright{}, attached now to an hour three hours away from the one it was named for. The word stayed put; the clock moved under it.
\end{cousinweb}

\subsection*{Your turn}
Say these aloud:

\begin{itemize}
\raggedright
  \item \emph{nōnus, nōna, nōnum}
  \item the squeeze — \emph{novem $\to$ nōnus}
  \item \emph{nōna hōra} — the ninth hour, once mid-afternoon
\end{itemize}

\begin{tcolorbox}[breakable,colback=teal!4,colframe=teal!35!black,title={Before you move on}]
How is \emph{nōnus} made out of \emph{novem}? (\textbf{By squeezing it shut}, not by adding — \emph{noven-} to \emph{nōn-}.) Which English word is \emph{nōna hōra}? (\emph{\textbf{Noon}}.) And what moved, the word or the hour? (\textbf{The hour} — the name stayed and the time it named drifted to midday.)
\end{tcolorbox}
\section[decimus, decima, decimum]{decimus — the tenth, and the ten complete}

\label{lesson:LA-C56-decimus}



\begin{quote}
Recalls first, then the tenth, which closes the series.

\begin{itemize}
\raggedright
  \item \emph{Recall:} say \emph{nōnus} — \textbf{R1}, one lesson back, before it decays
  \item \emph{Recall:} say \emph{quīntus} — \textbf{R2}, the first real retrieval, five to fifteen lessons back
  \item \emph{Recall:} say \emph{focus} — \textbf{R3}, durable, twenty to sixty lessons back
  \item \emph{Recall:} say \emph{tuus} — \textbf{R3}, a second word from the same distance
  \item \emph{Recall:} say \emph{pater and mater} — \textbf{R4}, recognition at distance, eighty lessons back or more
  \item \emph{Recall:} say \emph{frater and soror} — \textbf{R4}, a second word from the same distance
\end{itemize}
\end{quote}

\begin{cousinweb}
\textbf{decimus, decima, decimum} — \textquotedblleft{}\textbf{tenth}.\textquotedblright{}

\emph{Decem} is \textbf{ten}, and the ending goes on with the same worn-down middle you met in \emph{septimus}: \emph{decem} $\to$ \emph{decimus}. Remember that \emph{c} is hard — \emph{DEK-imus}.

\textbf{December} is this number in a name, and it is the last of the four the calendar froze.

\emph{Decima} on its own means \textquotedblleft{}a tenth part\textquotedblright{}, and that noun did a great deal of work. \textbf{Decimate} is the Roman army's punishment: a unit that broke was made to draw lots and one man in ten was killed by the other nine. A \textbf{dime} is a tenth of a dollar. A \textbf{tithe} is a tenth of a harvest, and \emph{décima} is still the word for it in Spanish.

That is the ten. Say them in order and you have a series that no cardinal count gives you: the place a thing holds, not how many of it there are.
\end{cousinweb}

\subsection*{Your turn}
Say these aloud:

\begin{itemize}
\raggedright
  \item \emph{decimus, decima, decimum} — \emph{DEK-imus}
  \item the join — \emph{decem $\to$ decimus}
  \item all ten in order, \emph{prīmus} to \emph{decimus}
\end{itemize}

\begin{tcolorbox}[breakable,colback=teal!4,colframe=teal!35!black,title={Before you move on}]
How is \emph{c} said in \emph{decimus}? (\textbf{Hard} — \emph{DEK-imus}.) What did the Roman army actually do when it decimated a unit? (\textbf{Killed one man in ten}, chosen by lot.) And how many of the ten ordinals have their own number visible inside them? (\textbf{Eight} — every one but \emph{prīmus} and \emph{secundus}.)
\end{tcolorbox}
